\chapter{Transport Layer}
\label{ch:transport}

\section*{Unit map: where the marks are}
The transport layer is the \textbf{heaviest midsem unit}: transport questions carried 20.5 of 27.5 (2022), 15.5 of 25 (2024) and 14.5 of 27.5 (2025) printed marks (from the PYQ bank).
Three question families appeared in all three papers: \keyword{ARQ loss-counting / utilisation numericals} (stop-and-wait, Go-Back-N),
\keyword{header decoding} (UDP and TCP hex dumps, urgent pointer) and \keyword{TCP connection management} (handshake, sequence numbers, states).
Read the table top-down; it is sorted by the chance $P$ that the topic appears in your paper.

\begin{table}[H]\centering\footnotesize\zebra
\begin{tabularx}{\linewidth}{L{48mm}cY}\toprule
\thead{Topic (section)} & \thead{P} & \thead{Must-know items}\\\midrule
Stop-and-wait (\S\ref{sec:T14}) & \prob{84} & window 1, seq.\ 0/1, timer; lost packet vs lost ACK; BDP; $U=\frac{1}{1+2a}$; loss counting\\
Go-Back-N (\S\ref{sec:T15}) & \prob{84} & $S_f,S_n,R_n$; $S_{size}\le 2^m-1$; cumulative ACK; ``every $k$-th lost'' count\\
UDP (\S\ref{sec:T17}) & \prob{80} & 8-byte header, hex decode, checksum = 0 means none; TCP vs UDP; DNS over UDP/TCP\\
TCP segment format (\S\ref{sec:T19}) & \prob{80} & header diagram, 6 flags, HLEN$\times4$, urgent pointer, hex decode\\
TCP connection management (\S\ref{sec:T20}) & \prob{59} & 3-way handshake with seq/ack, piggybacking, FIN, half-close, full trace\\
TCP congestion control (\S\ref{sec:T23}) & \prob{59} & slow start, additive increase, fast recovery, Tahoe vs Reno, $0.75W_{max}/RTT$\\
TCP windows and flow control (\S\ref{sec:T22}) & \prob{41} & rwnd = buffer $-$ waiting; window $=\min(rwnd,cwnd)$; window vs MSS\\
Selective Repeat (\S\ref{sec:T16}) & \prob{40} & $2^{m-1}$ windows, individual ACK, window-size derivation with failure diagram\\
TCP byte numbering (\S\ref{sec:T18}) & \prob{40} & seq.\ no.\ = first byte; SYN/FIN consume 1, pure ACK 0\\
TCP state diagram (\S\ref{sec:T21}) & \prob{35} & 11 states, client and server paths, TIME-WAIT = 2MSL\\
Transport services, ports, mux/demux (\S\ref{sec:T12}) & \prob{26} & process-to-process, port ranges, socket address\\
Flow/error control, simple protocol (\S\ref{sec:T13}) & \prob{14} & push vs pull, connectionless vs connection-oriented\\
Quality of Service (\S\ref{sec:T24}) & \prob{8} & reliability, delay, jitter, bandwidth; scheduling, shaping\\\bottomrule
\end{tabularx}
\caption{Transport layer at a glance. Lectures 16--21 of the notes \src{N_Network_Systems_compressed p60-90}.}
\end{table}


% =====================================================================
\section{Transport-layer services: process-to-process delivery, ports, multiplexing}\label{sec:T12}
\prob{26} \tier{C} \asked{2022 Q1(b), 2022 Q4(b)}

\begin{intuitionbox}
IP delivers to the host (the building); the \keyword{port number} picks the process (the flat). Many processes into one stream is \keyword{multiplexing}; one stream to the right process is \keyword{demultiplexing}.
\end{intuitionbox}

\begin{definitionbox}[Transport layer]
The transport layer provides \keyword{process-to-process communication} between two application layers. The communication
uses a \keyword{logical connection}: the two application layers assume an imaginary direct connection through which they send and
receive messages. \src{N_Network_Systems_compressed p60}
\end{definitionbox}

\begin{itemize}
  \item \keyword{Port number}: 16-bit integer, 0--65535, identifies a process on a host. A client uses a temporary
        (\emph{ephemeral}) port; a server uses a well-known port.
  \item \keyword{Socket address} = IP address + port number. A transport connection is identified by the pair of socket addresses.
  \item \keyword{Encapsulation}: header (ports, length, checksum) + message $=$ packet (UDP \emph{user datagram}, TCP \emph{segment}); \keyword{decapsulation} removes it at the receiver. \src{N_Network_Systems_compressed p60}
\end{itemize}

\begin{figure}[H]\centering
\begin{tikzpicture}[x=1cm,y=1cm]
  \draw[line width=5pt,color=ok!60] (0,0) -- (3,0);
  \draw[line width=5pt,color=info!60] (3.1,0) -- (8.9,0);
  \draw[line width=5pt,color=hl!60] (9,0) -- (13,0);
  \foreach \x/\t in {0/0,3/1023,3.1/1024,8.9/49151,9/49152,13/65535}{\draw[sgslate] (\x,0.2) -- (\x,-0.2);}
  \node[note,below left] at (3,-0.2) {1023}; \node[note,below right] at (3.1,-0.2) {1024};
  \node[note,below left] at (8.9,-0.2) {49151}; \node[note,below right] at (9,-0.2) {49152};
  \node[note,below] at (0,-0.2) {0}; \node[note,below] at (13,-0.2) {65535};
  \node[font=\scriptsize\bfseries,above] at (1.5,0.2) {Well-known};
  \node[font=\scriptsize\bfseries,above] at (6,0.2) {Registered};
  \node[font=\scriptsize\bfseries,above] at (11,0.2) {Dynamic / private};
  \node[note] at (1.5,-0.85) {assigned by ICANN};
  \node[note] at (6,-0.85) {registered with ICANN, not assigned};
  \node[note] at (11,-0.85) {temporary (client) ports};
\end{tikzpicture}
\caption{ICANN ranges of port numbers (not to scale). \src{N_Network_Systems_compressed p60}}
\end{figure}

\begin{figure}[H]\centering
\begin{tikzpicture}[x=1cm,y=1cm]
  \foreach \y/\n/\p in {1.4/P1/5000,0/P2/5001,-1.4/P3/5002}{
    \node[blk,minimum width=11mm] (s\n) at (0,\y) {\n};
    \node[blk,minimum width=11mm] (r\n) at (11,\y) {\n};
  }
  \node[blkhl,minimum height=36mm,minimum width=14mm] (mux) at (2.9,0) {MUX\\[1mm]\scriptsize transport\\\scriptsize layer};
  \node[blk,minimum height=10mm,minimum width=22mm] (ip) at (5.5,0) {IP / network};
  \node[blkhl,minimum height=36mm,minimum width=14mm] (dmx) at (8.1,0) {DEMUX\\[1mm]\scriptsize transport\\\scriptsize layer};
  \draw[arr] (sP1) -- node[elabel,above,sloped]{src 5000} (mux.west|-sP1);
  \draw[arr] (sP2) -- node[elabel,above]{src 5001} (mux.west|-sP2);
  \draw[arr] (sP3) -- node[elabel,below,sloped]{src 5002} (mux.west|-sP3);
  \draw[arrhl] (mux) -- (ip); \draw[arrhl] (ip) -- (dmx);
  \draw[arr] (dmx.east|-rP1) -- node[elabel,above]{dst 5000} (rP1);
  \draw[arr] (dmx.east|-rP2) -- node[elabel,above]{dst 5001} (rP2);
  \draw[arr] (dmx.east|-rP3) -- node[elabel,below]{dst 5002} (rP3);
  \node[note] at (0,-2.4) {sending host: many-to-one};
  \node[note] at (11,-2.4) {receiving host: one-to-many};
\end{tikzpicture}
\caption{Multiplexing at the sender (many processes, one transport entity) and demultiplexing at the receiver (destination port selects the process).}
\end{figure}

\begin{itemize}
  \item \textbf{Multiplexing} (sender, many-to-one): accept messages from many processes, add a header with source and destination ports, hand all to IP.
  \item \textbf{Demultiplexing} (receiver, one-to-many): read the \emph{destination port} and deliver each payload to the matching process.
\end{itemize}

\begin{keyconceptbox}[Minimum services over IP (2022 Q1b)]
Full list: process-to-process delivery (ports), encapsulation, mux/demux, flow, error and congestion control. The \emph{minimum} a transport protocol must add over IP is \keyword{process-to-process delivery} (mux/demux by port) plus \keyword{error detection} (checksum): exactly what UDP gives.
\end{keyconceptbox}

\begin{warningbox}
Host-to-host delivery (IP address) is the network layer; process-to-process (port) is the transport layer. Well-known ends at 1023, registered at 49151.
\end{warningbox}

\begin{answerplan}[2.5 marks, short note ``multiplexing vs demultiplexing'']
\stepno{1} Define the transport layer as process-to-process delivery \quad \stepno{2} Draw the MUX/DEMUX figure with port labels \quad
\stepno{3} Give the comparison table: where, many-to-one vs one-to-many, which port field is used \quad \stepno{4} Keywords: \keyword{socket address}, \keyword{port number}, \keyword{logical connection}.
Asked in 2022 Q1(b) and Q4(b); solved on p.\,\pageref{q:MS2022-1b} and p.\,\pageref{q:MS2022-4b}.
\par\smallskip{\color{sgteal}\textbf{Self-check.}} \textbf{T12.1} Which ICANN range does port 52100 fall in, and who normally uses such a port?\quad
\textbf{T12.2} Which header field does the receiver use to demultiplex?
\end{answerplan}

% =====================================================================
\section{Flow control, error control, connectionless vs connection-oriented service, simple protocol}\label{sec:T13}
\prob{14} \tier{D} \asked{not yet asked; seniors' list: ``flow control in the transport layer''}

\begin{intuitionbox}
A fast talker needs a note-taker who can say ``wait'' (\keyword{flow control}) and ask for lost pages again (\keyword{error control}). Phoning first is connection-oriented; posting postcards is connectionless.
\end{intuitionbox}

\begin{definitionbox}[Flow control and error control]
\textbf{Flow control} balances the producer's rate against the consumer's. \emph{Pushing} (producer sends when ready) needs it; \emph{pulling} (consumer requests) does not. \textbf{Error control} makes delivery reliable: discard corrupted packets, resend lost or discarded ones, discard duplicates, buffer out-of-order packets until the gaps fill. \src{N_Network_Systems_compressed p61-62}
\end{definitionbox}

\begin{figure}[H]\centering
\begin{tikzpicture}[x=1cm,y=1cm]
  \node[blk,minimum width=26mm] (sa) at (0,2.6) {Producer\\\scriptsize sender application};
  \node[blk,minimum width=26mm] (st) at (0,0) {Consumer + producer\\\scriptsize sender transport};
  \node[blk,minimum width=26mm] (rt) at (8,0) {Producer + consumer\\\scriptsize receiver transport};
  \node[blk,minimum width=26mm] (ra) at (8,2.6) {Consumer\\\scriptsize receiver application};
  \draw[arr] ([xshift=-5mm]sa.south) -- node[elabel,left=1pt]{messages pushed} ([xshift=-5mm]st.north);
  \draw[arrok,dashed] ([xshift=5mm]st.north) -- node[elabel,right=1pt]{flow control} ([xshift=5mm]sa.south);
  \draw[arrhl] (st) -- node[elabel,above]{packets pushed} (rt);
  \draw[arr] ([xshift=5mm]rt.north) -- node[elabel,right=1pt]{messages pulled} ([xshift=5mm]ra.south);
  \draw[arr,dashed] ([xshift=-5mm]ra.south) -- node[elabel,left=1pt]{requests} ([xshift=-5mm]rt.north);
  \draw[arrok,dashed] (rt.south) -- ++(0,-0.8) -| node[elabel,pos=0.25]{flow control (buffer full / vacancy)} (st.south);
  \draw[arrbad] ([yshift=-3mm]rt.west) -- node[elabel,below]{error control (ACK / resend)} ([yshift=-3mm]st.east);
\end{tikzpicture}
\caption{Flow control at the transport layer: two feedback loops, receiver transport $\to$ sender transport and sender transport $\to$ sender application. \src{N_Network_Systems_compressed p61-62}}
\end{figure}

\textbf{How it is done.} Buffers implement flow control: when the receiving buffer is full the receiver tells the sender to stop; when
a slot frees up it says ``resume''. Error control uses \keyword{sequence numbers} (a packet has an $m$-bit number, so numbers run $0$ to $2^m-1$
and wrap around), \keyword{acknowledgments} (ACK) and \keyword{timers}. Flow and error control are combined in one mechanism: the \keyword{sliding window}.



\textbf{Connectionless} (UDP): no set-up, packets independent and unnumbered, may arrive late, lost or out of order, no flow/error control.
\textbf{Connection-oriented} (TCP): establishment, data transfer, teardown; numbered bytes; in-order delivery; flow and error control.

\begin{keyconceptbox}[Simple protocol]
The \keyword{simple protocol} has \textbf{no flow control and no error control}. It assumes the network never loses or corrupts a packet
and the receiver can process every packet immediately, so there are no ACKs and no sequence numbers. The sender transport sends a packet
whenever its application makes a \emph{request}; the receiver transport delivers each packet as it \emph{arrives}. Both sides stay in a single
\textsf{ready} state. \src{N_Network_Systems_compressed p64}
\end{keyconceptbox}

\begin{warningbox}
Flow control protects the \emph{receiver} (end-to-end); congestion control (\S\ref{sec:T23}) protects the \emph{network}. Do not mix them.
\end{warningbox}

\begin{answerplan}[2.5 marks, ``flow control at the transport layer'']
Define flow control (producer/consumer rates), contrast push and pull, draw the four-box figure with the two flow-control loops, mention buffers
(stop/resume). For TCP add: the receiver advertises \keyword{rwnd} in every segment, sender window $=\min(rwnd,cwnd)$ (\S\ref{sec:T22}).
\par\smallskip{\color{sgteal}\textbf{Self-check.}} \textbf{T13.1} Why does pulling not need flow control?\quad \textbf{T13.2} Name the four jobs of transport-layer error control.\quad
\textbf{T13.3} Why can the simple protocol work without sequence numbers?
\end{answerplan}

% =====================================================================
\section{Stop-and-wait protocol}\label{sec:T14}
\prob{84} \tier{B} \asked{2022 Q3(a), Q3(b) [2nd Q3]; 2024 Q1(c); 2025 Q4(a)}

\begin{intuitionbox}
Pass one bucket, wait for ``got it'', pass the next: safe, but the link is mostly idle. Every utilisation numerical measures how idle.
\end{intuitionbox}

\begin{definitionbox}[Stop-and-wait protocol]
A connection-oriented protocol that uses both flow and error control. Both sender and receiver use a sliding window of
\textbf{size 1}: the sender sends one packet and waits for its acknowledgment before sending the next one. \src{N_Network_Systems_compressed p65}
\end{definitionbox}

\begin{algorithmbox}[Stop-and-wait ($S$ = sender's next number, $R$ = receiver's expected number)]
\begin{enumerate}
  \item Sequence numbers modulo 2; the \keyword{ACK number} is the \emph{next} packet expected.
  \item Sender sends packet $S$, keeps a copy, \keyword{starts the timer}; on \textsf{ACK} $=S+1$ it stops the timer and slides ($S\leftarrow S+1$).
  \item Receiver gets error-free packet $R$: delivers it, $R\leftarrow R+1 \pmod 2$, sends \textsf{ACK}~$R$. A duplicate ($\neq R$) is \textbf{discarded} but \textsf{ACK}~$R$ is \textbf{resent}; a corrupted packet is silently dropped.
  \item \textbf{Timeout} (packet or ACK lost): resend the same packet, restart the timer.
\end{enumerate}
\end{algorithmbox}

\stepanim{4}{2}{0.48\linewidth}{Stop-and-wait flow diagram, one event per frame (time runs downward): normal exchange, lost packet, lost ACK, duplicate discarded. \src{N_Network_Systems_compressed p65}}{%
  \path[use as bounding box] (-1.3,-3.95) rectangle (6.5,0.85);
  \node[blk,minimum width=13mm,minimum height=5mm,font=\scriptsize,inner sep=1pt] at (0,0.15) {Sender};
  \node[blk,minimum width=13mm,minimum height=5mm,font=\scriptsize,inner sep=1pt] at (5.2,0.15) {Receiver};
  \draw[lifeline] (0,-0.15) -- (0,-3.35); \draw[lifeline] (5.2,-0.15) -- (5.2,-3.35);
  \ifcase\f\or
    \draw[msg] (0,-0.4) -- node[elabel,font=\scriptsize]{Pkt 0} (5.2,-1.1);
    \node[note,anchor=west] at (5.3,-1.1) {$R:0\to1$};
    \draw[msgok] (5.2,-1.3) -- node[elabel,font=\scriptsize]{ACK 1} (0,-2.0);
    \node[note,anchor=east] at (-0.1,-2.0) {$S:0\to1$};
    \node[note,text width=76mm] at (2.6,-3.7) {Pkt 0 arrives, ACK 1 (next expected) returns; sender slides.};
  \or
    \draw[msglost] (0,-0.4) -- node[elabel,font=\scriptsize]{Pkt 1} (2.6,-0.75);
    \node[text=bad,font=\bfseries] at (2.72,-0.77) {$\times$};
    \draw[timeout] (-0.15,-0.4) -- node[left=4pt,note,text=bad]{T-out} (-0.15,-1.6);
    \draw[msg] (0,-1.7) -- node[elabel,font=\scriptsize]{Pkt 1 (resent)} (5.2,-2.4);
    \draw[msgok] (5.2,-2.55) -- node[elabel,font=\scriptsize]{ACK 0} (0,-3.2);
    \node[note,text width=76mm] at (2.6,-3.7) {Pkt 1 lost: no ACK, timer expires, Pkt 1 resent.};
  \or
    \draw[msg] (0,-0.4) -- node[elabel,font=\scriptsize]{Pkt 0} (5.2,-1.1);
    \node[note,anchor=west] at (5.3,-1.1) {$R:0\to1$};
    \draw[msglost] (5.2,-1.3) -- node[elabel,font=\scriptsize]{ACK 1} (2.6,-1.65);
    \node[text=bad,font=\bfseries] at (2.48,-1.67) {$\times$};
    \node[note,text width=76mm] at (2.6,-3.7) {Pkt 0 delivered, but ACK 1 is lost.};
  \or
    \draw[timeout] (-0.15,-0.4) -- node[left=4pt,note,text=bad]{T-out} (-0.15,-1.6);
    \draw[msg] (0,-1.7) -- node[elabel,font=\scriptsize]{Pkt 0 (resent)} (5.2,-2.4);
    \node[note,anchor=west,text=hl,align=left] at (5.3,-2.4) {dup:\\discard};
    \draw[msgok] (5.2,-2.55) -- node[elabel,font=\scriptsize]{ACK 1} (0,-3.2);
    \node[note,text width=76mm] at (2.6,-3.7) {Timeout: receiver discards the duplicate, resends ACK 1.};
  \fi
}

\begin{formulabox}[Stop-and-wait performance (notes p91)]
\[
T_t=\frac{L}{B},\quad T_p=\frac{d}{v},\quad a=\frac{T_p}{T_t},\qquad
U_{SW}=\frac{T_t}{T_t+2T_p}=\frac{1}{1+2a},\qquad U_{W}=\min\!\left(1,\frac{W}{1+2a}\right)
\]
$B$ = bandwidth (bps), $L$ = packet length (bits), $d$ = distance, $v$ = propagation speed, $W$ = packets sent before waiting.
\textbf{Intuition:} $\text{BDP}=B\times RTT$ bits fill the round trip; stop-and-wait puts only $L$ of them in the pipe, so $U\approx L/\text{BDP}$ (the textbook short form, which ignores $T_t$). \src{N_Network_Systems_compressed p91}
\end{formulabox}

\needspace{8\baselineskip}
\begin{question}{Worked example 1 -- Variant of MS2025-4a}{3}[\tier{B}]
Stop-and-wait ARQ, line bandwidth 2 Mbps, 1 bit takes 10 ms for a round trip, data packets of 2000 bits. Find the bandwidth-delay
product and the link utilisation; then find the utilisation if up to 8 packets may be sent before waiting for ACKs.
\end{question}
\begin{solution}[BDP and utilisation (check: \texttt{checks/T14\_ex\_sw\_util.py})]
$\text{BDP}=2\times10^{6}\times10\times10^{-3}=\answer{$20\,000$ bits}$.\quad
$T_t=\frac{2000}{2\times10^6}=1$ ms, $T_p=\frac{10}{2}=5$ ms, $a=5$.\\
$U_{SW}=\frac{1}{1+2a}=\frac{1}{11}=\answer{$9.09\%$}$,\quad $U_{W=8}=\frac{8}{11}=\answer{$72.7\%$}$.
\emph{Short form} $L/\text{BDP}$: 10\% and 80\% (the textbook answer to this PYQ style); write one line saying which you use.
\end{solution}

\needspace{8\baselineskip}
\begin{question}{Worked example 2 -- Variant of MS2024-1c}{1}[\tier{B}]
Stop-and-wait: every 3rd packet transmitted is lost and 10 packets must be delivered. How many transmissions are needed?
\end{question}
\begin{solution}[count transmissions (check: \texttt{checks/T14\_ex\_sw\_loss.py})]
Transmissions 3, 6, 9, 12 are lost (retransmissions count). Need $T-\lfloor T/3\rfloor=10$ with $T$ not lost: $T=N+\left\lfloor\frac{N-1}{k-1}\right\rfloor=10+4=\answer{14}$.
\end{solution}

\needspace{8\baselineskip}
\begin{question}{Worked example 3 -- Expected (seniors' list)}{2.5}[\tier{B}]
Earth to a distant planet $\approx 9\times10^{10}$ m. Stop-and-wait on a 64 Mbps link, frame 32 KB, $v=3\times10^{8}$ m/s. Find the channel utilisation.
\end{question}
\begin{solution}[very large $a$ (check: \texttt{checks/T14\_ex\_planet.py})]
1 KB $=1024$ B (convention): $T_t=\frac{32\times1024\times8}{64\times10^{6}}=4.096$ ms, $T_p=\frac{9\times10^{10}}{3\times10^{8}}=300$ s, $a=73\,242$,
$U=\frac{1}{1+2a}\approx\answer{$6.83\times10^{-6}$ ($0.00068\%$)}$. (1 KB $=1000$ B: $6.67\times10^{-6}$.)
\end{solution}

\begin{warningbox}
``Every $k$-th lost'' counts \emph{transmissions} (retransmissions included) and the last transmission cannot be a lost one. ``1 bit takes 20 ms for a round trip'' already \emph{is} the RTT. ACK $=$ next expected (Pkt 0 $\to$ ACK 1).
\end{warningbox}

\begin{answerplan}[2.5 marks, ``flow diagram of stop-and-wait with packet lost and ACK lost'']
\stepno{1} Two-line definition (window size 1, wait for ACK, timer) \quad \stepno{2} Draw the 4-frame diagram above as one timeline: normal, lost packet + timeout, lost ACK + duplicate discarded \quad
\stepno{3} Keywords: \keyword{sequence number mod 2}, \keyword{ACK = next expected}, \keyword{timeout}, \keyword{duplicate discarded}. Asked in 2022 (2nd Q3a); solved on p.\,\pageref{q:MS2022-3Ba}.
Numericals: 2024 Q1(c) on p.\,\pageref{q:MS2024-1c}; 2022 (2nd Q3b) on p.\,\pageref{q:MS2022-3Bb}; 2025 Q4(a) on p.\,\pageref{q:MS2025-4a}.
\par\smallskip{\color{sgteal}\textbf{Self-check.}} \textbf{T14.1} Why are sequence numbers needed at all if only one packet is outstanding?\quad
\textbf{T14.2} $B=1$ Mbps, $RTT=20$ ms, $L=1000$ bits: what is $a$ and the exact utilisation?\quad
\textbf{T14.3} Every 5th transmission is lost; 8 packets. How many transmissions?
\end{answerplan}

% =====================================================================
\section{Go-Back-N protocol}\label{sec:T15}
\prob{84} \tier{A} \asked{2022 Q1(a); 2024 Q4(b); 2025 Q3(b) (the same numerical three times)}

\begin{intuitionbox}
Keep many packets in flight. The lazy receiver accepts only the packet it expects, so one loss forces a resend of that packet and everything sent after it.
\end{intuitionbox}

\begin{definitionbox}[Go-Back-N (GBN)]
Used to overcome the limitation of stop-and-wait: \textbf{multiple packets are in transit} while the sender waits for acknowledgment,
so the bandwidth is fully used. A \keyword{cumulative acknowledgment} is used, and a sliding window is kept at both sender and receiver.
Sequence numbers are modulo $2^m$ (for $m=3$: 0 to 7). \src{N_Network_Systems_compressed p66}
\end{definitionbox}

\begin{figure}[H]\centering
\begin{tikzpicture}[x=1cm,y=1cm]
  \foreach \v/\s [count=\i] in {6/celldim,7/celldim,0/cellhl,1/cellhl,2/cellhl,3/cellhl,4/cellok,5/cellok,6/cellok,7/cell,0/cell}{\node[\s] at (\i*0.75,0) {\v};}
  \draw[hl,line width=1.3pt] (1.875,-0.38) rectangle (7.125,0.38);
  \draw[arr] (2.25,1.15) node[above,font=\scriptsize]{$S_f$ first outstanding} -- (2.25,0.42);
  \draw[arr] (5.25,1.15) node[above,font=\scriptsize]{$S_n$ next to send} -- (5.25,0.42);
  \draw[decorate,decoration={brace,mirror,amplitude=4pt}] (0.4,-0.5) -- node[below=4pt,note]{ACKed,\\purged} (1.85,-0.5);
  \draw[decorate,decoration={brace,mirror,amplitude=4pt}] (1.9,-0.5) -- node[below=4pt,note]{outstanding\\(sent, no ACK)} (4.85,-0.5);
  \draw[decorate,decoration={brace,mirror,amplitude=4pt}] (4.9,-0.5) -- node[below=4pt,note]{can be\\sent} (7.1,-0.5);
  \draw[decorate,decoration={brace,mirror,amplitude=4pt}] (7.15,-0.5) -- node[below=4pt,note]{cannot be\\sent} (8.6,-0.5);
  \draw[<->,sgslate] (1.875,-1.75) -- node[elabel]{$S_{size}=2^m-1=7$} (7.125,-1.75);
  % receive window
  \begin{scope}[xshift=10.2cm]
    \foreach \v/\s [count=\i] in {3/celldim,4/celldim,5/cellhl,6/cell,7/cell}{\node[\s] at (\i*0.75,0) {\v};}
    \draw[arr] (2.25,1.15) node[above,font=\scriptsize]{$R_n$ next expected} -- (2.25,0.42);
    \node[note] at (1.1,-0.85) {received,\\ACKed};
    \node[note] at (3.4,-0.85) {cannot be\\received};
    \node[elabel] at (2.25,-1.75) {$R_{size}=1$};
  \end{scope}
\end{tikzpicture}
\caption{GBN send window ($m=3$, left) and receive window (right). \src{N_Network_Systems_compressed p66-67}}
\end{figure}

\begin{algorithmbox}[Go-Back-N]
\begin{enumerate}
  \item Sender may have up to $S_{size}=2^m-1$ outstanding packets, from $S_f$ to $S_n-1$. One \keyword{timer}, for the first outstanding packet $S_f$.
  \item Receiver window has size 1 ($R_n$). A packet numbered $R_n$ is delivered, $R_n$ advances and an ACK with ackNo $=R_n$ (next expected) is sent.
        Any other packet (out of order or duplicate) is \textbf{discarded}; the receiver only re-acknowledges.
  \item ACKs are \keyword{cumulative}: ACK $n$ confirms every packet before $n$, so $S_f\leftarrow n$. Example: window $4\ldots2$ with $S_f=4$; ACK 6 arrives $\Rightarrow S_f=6$ and packets 4, 5 are purged.
  \item \textbf{Timeout}: resend \emph{all} outstanding packets, $S_f$ to $S_n-1$ (``go back N'').
\end{enumerate}
\end{algorithmbox}

\textbf{Method for ``every $k$-th lost'' (GBN):} number the transmissions; after a lost packet $p$ the sender still sends $p+1\ldots p+W-1$ (discarded), then times out and restarts at $p$; count until the last packet is delivered.

\stepanim{3}{1}{\linewidth}{GBN transmission strip for the worked example (10 packets, $W=4$, every 6th transmission lost). Top row = transmission number; green = accepted, red = lost, grey = arrives out of order and is discarded.}{%
  \path[use as bounding box] (-1.5,-1.45) rectangle (11.3,0.95);
  \node[font=\scriptsize\bfseries,anchor=east] at (0.2,0) {pkt};
  \node[font=\scriptsize,anchor=east,text=muted] at (0.2,0.52) {tx};
  \foreach \p/\s/\r [count=\i] in {1/cellok/1,2/cellok/1,3/cellok/1,4/cellok/1,5/cellok/1,6/cellbad/1,7/celldim/1,8/celldim/1,9/celldim/1,6/cellok/2,7/cellok/2,8/cellbad/2,9/celldim/2,10/celldim/2,8/cellok/3,9/cellok/3,10/cellok/3}{
    \ifnum\r>\f
      \node[cell,minimum width=5.6mm,minimum height=5.6mm,draw=muted!40] at (\i*0.62,0) {};
    \else
      \node[\s,minimum width=5.6mm,minimum height=5.6mm,font=\scriptsize] at (\i*0.62,0) {\p};
    \fi
    \node[font=\scriptsize,text=muted] at (\i*0.62,0.52) {\i};
  }
  \draw[decorate,decoration={brace,mirror,amplitude=4pt}] (0.33,-0.38) -- node[below=4pt,note]{round 1} (5.85,-0.38);
  \ifnum\f>1 \draw[decorate,decoration={brace,mirror,amplitude=4pt}] (5.95,-0.38) -- node[below=4pt,note]{round 2} (8.97,-0.38); \fi
  \ifnum\f>2 \draw[decorate,decoration={brace,mirror,amplitude=4pt}] (9.07,-0.38) -- node[below=4pt,note]{round 3} (10.85,-0.38); \fi
  \ifcase\f\or
    \node[note,anchor=west] at (-1.4,-1.2) {tx 6 (pkt 6) lost; window 6--9 already sent, so 7, 8, 9 are discarded; timeout on 6.};
  \or
    \node[note,anchor=west] at (-1.4,-1.2) {go back to 6: 6, 7 accepted; tx 12 (pkt 8) lost; 9, 10 discarded; timeout on 8.};
  \or
    \node[note,anchor=west,text=ok] at (-1.4,-1.2) {go back to 8: 8, 9, 10 accepted. Total = 17 transmissions.};
  \fi
}

\needspace{8\baselineskip}
\begin{question}{Worked example -- Variant of MS2024-4b}{2.5}[\tier{A}]
A sender must deliver 10 packets with Go-Back-N, window size 4. Every 6th packet that it transmits is lost; no ACK is ever lost. How many packets does it transmit?
\end{question}
\begin{solution}[transmission strip (check: \texttt{checks/T15\_ex\_gbn\_loss.py})]
\begin{tabular}{@{}ll@{}}
Round 1: & 1 2 3 4 5 \textcolor{bad}{\textbf{6}} 7 8 9 \quad (tx 1--9; 7--9 discarded)\\
Round 2: & 6 7 \textcolor{bad}{\textbf{8}} 9 10 \quad (tx 10--14; 9, 10 discarded)\\
Round 3: & 8 9 10 \quad (tx 15--17; tx 18 would be lost but nothing is left)
\end{tabular}\\[2pt]
Total transmissions $= 9+5+3=$ \answer{17}. The same simulation with 9 packets, $W=3$ and every 5th lost (the past-paper data) is the one solved in the PYQ chapter.
\end{solution}

\textbf{GBN numbers:} $S_{size}\le 2^m-1$, $R_{size}=1$, $U_{GBN}=\min\!\left(1,\frac{W}{1+2a}\right)$; bits needed for a window $W$: $m=\lceil\log_2(W+1)\rceil$.

\begin{warningbox}
After a lost packet the sender still sends the rest of its window ($p+1\ldots p+W-1$): wasted but \emph{counted}. The transmission counter never resets per round.
\end{warningbox}

\begin{answerplan}[2--2.5 marks, GBN loss count]
\stepno{1} State the rule (window $W$, cumulative ACK, receiver discards out-of-order, timeout resends all outstanding) \quad
\stepno{2} Draw the transmission strip or timing diagram (the 2025 paper demands a ``transition diagram'') \quad
\stepno{3} Count and box the answer. Asked in 2022 Q1(a), 2024 Q4(b), 2025 Q3(b); solved on p.\,\pageref{q:MS2022-1a}, p.\,\pageref{q:MS2024-4b} and p.\,\pageref{q:MS2025-3b}.
\par\smallskip{\color{sgteal}\textbf{Self-check.}} \textbf{T15.1} $m=4$: what is the largest GBN send window?\quad
\textbf{T15.2} 6 packets, $W=3$, every 4th transmission lost: how many transmissions?\quad
\textbf{T15.3} What does ACK 5 tell a GBN sender?
\end{answerplan}

% =====================================================================
\needspace{14\baselineskip}
\section{Selective-Repeat protocol; window size vs sequence-number bits}\label{sec:T16}
\prob{40} \tier{B} \asked{2024 Q4(a)}

\begin{intuitionbox}
Let the receiver buffer early packets so only the lost ones are resent. The price: both windows shrink to half the number space.
\end{intuitionbox}

\begin{definitionbox}[Selective-Repeat (SR)]
The \keyword{Selective-Repeat} protocol resends only selective packets, \textbf{those that are actually lost}.
Send window $S_{size}=2^{m-1}$; receive window $R_{size}=2^{m-1}$. \src{N_Network_Systems_compressed p69-70}
\end{definitionbox}

\begin{figure}[H]\centering
\begin{tikzpicture}[x=1cm,y=1cm]
  \node[font=\scriptsize\bfseries,anchor=east] at (0.2,0) {Sender};
  \foreach \v/\s [count=\i] in {14/celldim,15/celldim,0/cellhl,1/cellhl,2/cell,3/cellhl,4/cellok,5/cellok,6/cellok,7/cellok,8/cell,9/cell}{\node[\s,minimum width=6.5mm] at (\i*0.68,0) {\v};}
  \draw[hl,line width=1.3pt] (1.70,-0.38) rectangle (7.14,0.38);
  \draw[arr] (2.04,1.05) node[above,font=\scriptsize]{$S_f$} -- (2.04,0.42);
  \draw[arr] (4.76,1.05) node[above,font=\scriptsize]{$S_n$} -- (4.76,0.42);
  \node[note,anchor=west,align=left] at (8.75,0) {orange = outstanding, not ACKed\\white in window = ACKed out of order\\green = can be sent};
  \draw[<->,sgslate] (1.70,-0.65) -- node[elabel]{$S_{size}=2^{m-1}=8$} (7.14,-0.65);
  \begin{scope}[yshift=-2.7cm]
  \node[font=\scriptsize\bfseries,anchor=east] at (0.2,0) {Receiver};
  \foreach \v/\s [count=\i] in {1/celldim,2/celldim,3/cell,4/cellok,5/cell,6/cellok,7/cell,8/cell,9/cellok,10/cell,11/cell,12/cell}{\node[\s,minimum width=6.5mm] at (\i*0.68,0) {\v};}
  \draw[info,line width=1.3pt] (1.70,-0.38) rectangle (7.14,0.38);
  \draw[arr] (2.04,1.05) node[above,font=\scriptsize]{$R_n$} -- (2.04,0.42);
  \node[note,anchor=west,align=left] at (8.75,0) {green = received out of order, stored\\all delivered when 3 arrives};
  \draw[<->,sgslate] (1.70,-0.65) -- node[elabel]{$R_{size}=2^{m-1}=8$} (7.14,-0.65);
  \end{scope}
\end{tikzpicture}
\caption{SR send and receive windows for $m=4$ (numbers 0--15). \src{N_Network_Systems_compressed p69-70}}
\end{figure}

\begin{algorithmbox}[Selective Repeat]
\begin{enumerate}
  \item Sender: up to $2^{m-1}$ outstanding packets; \keyword{one timer per packet}; on timeout resend \emph{only that} packet.
  \item Receiver: accepts any packet inside its window, \keyword{buffers} out-of-order ones, and sends an \keyword{individual ACK} carrying the number of the packet received.
  \item When the packet at $R_n$ arrives, it and all consecutive buffered packets are delivered and the window slides.
  \item ACK interpretation (sender sent 0--5 and gets ACK 3): GBN and stop-and-wait read ``0, 1, 2 received, 3 expected''; SR reads ``packet 3 received'', saying nothing about the others. \src{N_Network_Systems_compressed p70}
\end{enumerate}
\end{algorithmbox}

\subsection*{Derivation: window size vs sequence-number bits}
With $m$ bits there are $2^m$ distinct numbers. The danger is a \emph{retransmitted old} packet that the receiver mistakes for a \emph{new} one with the same number.
\begin{itemize}
  \item \textbf{GBN.} Let $S_{size}=2^m$. The sender sends $0,\ldots,2^m-1$; all arrive, so $R_n$ wraps to 0. If all ACKs are lost, the timer expires and old packet 0
        is resent. The receiver expects 0 and \textbf{wrongly accepts the duplicate as new data}. With $S_{size}=2^m-1$ the receiver, after
        $0,\ldots,2^m-2$, expects $2^m-1\neq0$, so the resent 0 is correctly discarded. Hence $S_{size}\le 2^m-1$.
  \item \textbf{SR.} The receiver accepts anything inside a window of $R_{size}$ numbers. After receiving a full send window, the receive window has moved
        $S_{size}$ places. A resent old packet is safely rejected only if the old send window and the new receive window do not overlap:
        $S_{size}+R_{size}\le 2^m$. Equal windows give $S_{size}=R_{size}=2^{m}/2=2^{m-1}$. \src{N_Network_Systems_compressed p68, p71}
\end{itemize}

\begin{figure}[!htbp]\centering
\begin{minipage}[t]{0.49\linewidth}\centering
\begin{tikzpicture}[x=1cm,y=1cm]
  \node[blk,font=\scriptsize,minimum width=12mm] at (0,0.3) {Sender}; \node[blk,font=\scriptsize,minimum width=12mm] at (3,0.3) {Receiver};
  \draw[lifeline] (0,-0.1) -- (0,-5.2); \draw[lifeline] (3,-0.1) -- (3,-5.2);
  \foreach \k/\y in {0/-0.4,1/-1.1,2/-1.8,3/-2.5}{
    \draw[msg] (0,\y) -- node[elabel,above,sloped,pos=0.3]{Pkt \k} (3,\y-0.45);
    \draw[msglost] (3,\y-0.5) -- (2.45,\y-0.66);
    \node[text=bad,font=\scriptsize\bfseries] at (2.37,\y-0.68) {$\times$};
  }
  \node[note,right] at (3.05,-3.0) {$R_n=0$};
  \draw[timeout] (-0.12,-0.4) -- node[left=4pt,note,text=bad]{T-out} (-0.12,-3.9);
  \draw[msg] (0,-4.0) -- node[elabel,above,sloped]{Pkt 0 (old)} (3,-4.5);
  \node[note,text=bad,align=left,anchor=west] at (3.05,-4.6) {accepted\\as new!};
\end{tikzpicture}
\par\small (a) GBN, $m=2$, $S_{size}=2^m=4$ (wrong)
\end{minipage}\hfill
\begin{minipage}[t]{0.49\linewidth}\centering
\begin{tikzpicture}[x=1cm,y=1cm]
  \node[blk,font=\scriptsize,minimum width=12mm] at (0,0.3) {Sender}; \node[blk,font=\scriptsize,minimum width=12mm] at (3,0.3) {Receiver};
  \draw[lifeline] (0,-0.1) -- (0,-5.2); \draw[lifeline] (3,-0.1) -- (3,-5.2);
  \foreach \k/\y in {0/-0.4,1/-1.1,2/-1.8}{
    \draw[msg] (0,\y) -- node[elabel,above,sloped,pos=0.3]{Pkt \k} (3,\y-0.45);
    \draw[msglost] (3,\y-0.5) -- (2.45,\y-0.66);
    \node[text=bad,font=\scriptsize\bfseries] at (2.37,\y-0.68) {$\times$};
  }
  \node[note,right,align=left] at (3.05,-2.6) {window\\$\{3,0,1\}$};
  \draw[timeout] (-0.12,-0.4) -- node[left=4pt,note,text=bad]{T-out} (-0.12,-3.9);
  \draw[msg] (0,-4.0) -- node[elabel,above,sloped]{Pkt 0 (old)} (3,-4.5);
  \node[note,text=bad,align=left,anchor=west] at (3.05,-4.6) {stored\\as new!};
\end{tikzpicture}
\par\small (b) SR, $m=2$, windows $=3>2^{m-1}$ (wrong)
\end{minipage}
\caption{What goes wrong when the window is too large (all ACKs lost). With $S_{size}=3$ in (a) the receiver would expect 3 and discard old 0; with windows of 2 in (b) its window would be $\{2,3\}$ and old 0 would be rejected. \src{N_Network_Systems_compressed p68, p71}}
\end{figure}

\needspace{8\baselineskip}
\begin{question}{Worked example -- Variant of MS2024-4a}{2.5}[\tier{B}]
Sequence numbers use 4 bits. Find the maximum send and receive windows for GBN and SR.
\end{question}
\begin{solution}[windows for $m=4$ (check: \texttt{checks/T16\_ex\_window.py})]
Numbers $0$--$15$. GBN: $S_{size}=2^4-1=\answer{15}$, $R_{size}=1$.\quad SR: $S_{size}=R_{size}=2^{3}=\answer{8}$.
Then draw diagram (a) and (b) above with your $m$ to show the failure for a larger window.
\end{solution}

\needspace{8\baselineskip}
\begin{question}{Worked example -- Expected (cluster C27)}{1}[\tier{D}]
SR with $m=3$ (window 4). The sender has $S_f=2$, $S_n=5$. Which packets are outstanding and how many more can it send?
\end{question}
\begin{solution}[read the window (check: \texttt{checks/T16\_ex\_window.py})]
Outstanding: $S_f\ldots S_n-1=\{2,3,4\}$ (3 packets). Window 4 $\Rightarrow$ it may send \answer{1 more (packet 5)} before $S_f$ moves.
\end{solution}

\begin{table}[!htbp]\centering\footnotesize\zebra
\begin{tabularx}{\linewidth}{L{30mm}YYY}\toprule
\thead{Aspect} & \thead{Stop-and-wait} & \thead{Go-Back-N} & \thead{Selective Repeat}\\\midrule
Send window & 1 & $2^m-1$ & $2^{m-1}$\\
Receive window & 1 & 1 & $2^{m-1}$\\
ACK & next expected & cumulative, next expected & individual, number of packet received\\
Timers & 1 & 1 (oldest outstanding) & 1 per packet\\
On loss resend & that packet & all outstanding packets & only the lost packet\\
Out-of-order packet & discarded & discarded & buffered\\
Efficiency & $\frac{1}{1+2a}$ & $\frac{W}{1+2a}$ (wastes resends) & $\frac{W}{1+2a}$ (fewest resends)\\\bottomrule
\end{tabularx}
\caption{The three ARQ protocols side by side (a favourite 2.5-mark comparison).}
\end{table}

\begin{warningbox}
SR window is $2^{m-1}$, not $2^m-1$ ($m=3$: SR 4, GBN 7). SR ACK names the packet received; GBN ACK names the next one expected.
\end{warningbox}

\begin{answerplan}[2.5 marks, ``derive the relationship between window size and sequence number'']
\stepno{1} $m$ bits $\Rightarrow 2^m$ numbers \quad \stepno{2} GBN: show the $2^m$ failure diagram (all ACKs lost, old 0 accepted) and the $2^m-1$ fix \quad
\stepno{3} SR: overlap argument, $S_{size}+R_{size}\le2^m \Rightarrow 2^{m-1}$, with failure diagram \quad \stepno{4} Plug in the given $m$.
Asked in 2024 Q4(a); solved on p.\,\pageref{q:MS2024-4a}.
\par\smallskip{\color{sgteal}\textbf{Self-check.}} \textbf{T16.1} A GBN sender needs a window of 31 packets. How many sequence-number bits?\quad
\textbf{T16.2} Why does SR need one timer per packet but GBN only one?
\end{answerplan}

% =====================================================================
\section{UDP: services, applications, user-datagram format, checksum}\label{sec:T17}
\prob{80} \tier{B} \asked{2022 Q1(c), Q3(a); 2024 Q1(b); 2025 Q1(b)}

\begin{intuitionbox}
UDP = IP + ports + checksum: no handshake, ACKs, ordering or throttling. Tiny and fast, so right for DNS and live media, wrong for file transfer.
\end{intuitionbox}

\begin{definitionbox}[User Datagram Protocol]
UDP is a \keyword{connectionless, unreliable} transport protocol. Its packet, the \keyword{user datagram}, has a fixed
\textbf{8-byte header} followed by data; the whole datagram is 8 to 65\,535 bytes. Header fields (16 bits each): source port number,
destination port number, total length (header + data), checksum. \src{N_Network_Systems_compressed p72}
\end{definitionbox}

\begin{figure}[H]\centering
\begin{tikzpicture}[x=4mm,y=8mm]
  \draw[draw=sgnavy,fill=hl!15] (0,2.2) rectangle (8,1.5); \node[font=\scriptsize] at (4,1.85) {Header (8 B)};
  \draw[draw=sgnavy,fill=sgtint] (8,2.2) rectangle (32,1.5); \node[font=\scriptsize] at (20,1.85) {Data};
  \draw[<->,sgslate] (0,2.5) -- node[elabel]{8 to 65\,535 bytes} (32,2.5);
  \draw[muted,dashed] (0,1.5) -- (0,0); \draw[muted,dashed] (8,1.5) -- (32,0);
  \foreach \b in {0,16,31}{\node[note] at ({\b+0.5},0.25) {\b};}
  \foreach \x/\w/\r/\t in {0/16/0/{Source port number (16)},16/16/0/{Destination port number (16)},0/16/1/{Total length (16)},16/16/1/{Checksum (16)}}{
    \draw[draw=sgnavy,fill=sgtint,line width=0.6pt] (\x,-\r) rectangle ++(\w,-1);
    \node[font=\small] at ({\x+\w/2},{-\r-0.5}) {\t};
  }
\end{tikzpicture}
\caption{(a) UDP user datagram and (b) its 8-byte header format. \src{N_Network_Systems_compressed p72}}
\end{figure}

\needspace{8\baselineskip}
\begin{question}{Worked example -- from the notes (hex dump decoding)}{2}[\tier{B}]
A UDP header in hexadecimal is \texttt{CB84 000D 001C 001C}. Find the source port, destination port, total length, data length, the direction of
the packet and the client process. \src{N_Network_Systems_compressed p72-73}
\end{question}
\begin{solution}[split into four 16-bit (4-hex-digit) fields (check: \texttt{checks/T17\_ex\_udp.py})]
\begin{tabular}{@{}lll@{}}
Source port & \texttt{CB84}$_{16}$ & $=12\cdot4096+11\cdot256+8\cdot16+4=$ \answer{52100}\\
Destination port & \texttt{000D}$_{16}$ & $=$ \answer{13}\\
Total length & \texttt{001C}$_{16}$ & $=$ \answer{28 bytes}\\
Data length & $28-8$ & $=$ \answer{20 bytes}\\
Checksum & \texttt{001C} & non-zero, so the sender \emph{did} compute a checksum
\end{tabular}\\[2pt]
Destination port 13 is well-known (\textsf{daytime}) and 52100 is ephemeral, so the packet goes \textbf{client $\to$ server} and the client process is \textbf{Daytime}.
\end{solution}

\needspace{8\baselineskip}
\begin{question}{Worked example -- Variant of MS2022-1c}{1}[\tier{B}]
A client sends 40 bytes of data using UDP. What is the transmission efficiency at the UDP level?
\end{question}
\begin{solution}[efficiency $=\frac{\text{data}}{\text{data}+\text{header}}$ (check: \texttt{checks/T17\_ex\_udp.py})]
$\eta=\dfrac{40}{40+8}=\answer{$83.3\%$}$. (If the 20-byte IP header is also counted: $40/68=58.8\%$; the question says ``at the UDP level'', so use 8 bytes.)
\end{solution}

\textbf{Well-known ports} (``Table 3.1'' of the notes): 7 Echo (UDP), 9 Discard (UDP), 11 Users (UDP, TCP), 13 Daytime (UDP, TCP); textbook addition: 53 DNS (UDP, TCP). \src{N_Network_Systems_compressed p73}

\textbf{UDP services} \src{N_Network_Systems_compressed p73-74}: process-to-process communication (socket addresses); connectionless (independent, unnumbered datagrams); \textbf{no} flow control, \textbf{no} error control except the checksum (a corrupted datagram is silently discarded); no congestion control; encapsulation, queuing per port, multiplexing/demultiplexing. The \keyword{checksum} covers a \keyword{pseudoheader} (source IP, destination IP, 8 zero bits, protocol 17, UDP length), the header and the data padded to 16 bits.

\textbf{Three checksum cases} \src{N_Network_Systems_compressed p74-75}: (1) sender skips the checksum $\Rightarrow$ field \textbf{all 0s}; (2) the sum is all 1s, so the complement would be all 0s: it is complemented again and \textbf{all 1s} is sent; (3) a sum of all 0s \textbf{never happens}, because pseudoheader fields (IP addresses, protocol) are non-zero.

\textbf{UDP applications:} short request--response (DNS), real-time audio/video (a late packet is useless, so no error control is fine), SNMP, TFTP, RIP, DHCP, multicasting. \src{N_Network_Systems_compressed p75}

\begin{table}[H]\centering\footnotesize\zebra
\begin{tabularx}{\linewidth}{L{30mm}YY}\toprule
\thead{Aspect} & \thead{TCP} & \thead{UDP}\\\midrule
Connection & connection-oriented (3-way handshake) & connectionless\\
Reliability & reliable: ACK, retransmission, sequence numbers & unreliable, best effort\\
Ordering & in-order byte stream & datagrams may arrive out of order\\
Flow / congestion control & yes (rwnd, cwnd) & none\\
Header & 20--60 bytes & 8 bytes\\
Typical use & HTTP, SMTP, FTP, file transfer & DNS query, SNMP, TFTP, streaming, VoIP\\\bottomrule
\end{tabularx}
\caption{TCP vs UDP. For ``why is TCP preferred for data transfer'' quote rows 2--4.}
\end{table}

\begin{examtip}[DNS port 53: UDP \emph{or} TCP]
The notes write ``DNS uses TCP protocol with port number 53'' \src{N_Network_Systems_compressed p72}. DNS uses port 53 on \emph{both}.
\textbf{UDP} is the normal choice for queries: one short request and one short reply (up to 512 bytes), no connection set-up, the client simply retries.
\textbf{TCP} is used when the response is larger than 512 bytes (the UDP reply comes back truncated) and for \keyword{zone transfers} between primary and secondary servers, where reliability matters.
This is exactly 2022 Q3(a).
\end{examtip}

\begin{warningbox}
Length field = header + data (subtract 8). Checksum \texttt{0000} means ``not calculated''. Read the dump in 4-hex-digit fields.
\end{warningbox}

\begin{answerplan}[1--2.5 marks]
Hex dump: split into 4-digit fields, convert, subtract 8, comment on the checksum and on the well-known port. TCP vs UDP: give the table, then one line on the
application asked (data transfer needs reliability and order $\Rightarrow$ TCP). Asked in 2024 Q1(b), 2022 Q1(c), 2022 Q3(a), 2025 Q1(b); solved on
p.\,\pageref{q:MS2024-1b}, p.\,\pageref{q:MS2022-1c}, p.\,\pageref{q:MS2022-3a} and p.\,\pageref{q:MS2025-1b}.
\par\smallskip{\color{sgteal}\textbf{Self-check.}} \textbf{T17.1} A UDP header has total length \texttt{0040}$_{16}$. How many data bytes?\quad
\textbf{T17.2} What protocol number goes in the UDP pseudoheader?\quad
\textbf{T17.3} Give two reasons real-time video prefers UDP.
\end{answerplan}

% =====================================================================
\section{TCP services and features: byte numbering, sequence and acknowledgment numbers}\label{sec:T18}
\prob{40} \tier{B} \asked{2024 Q1(e), Q5(a)}

\begin{intuitionbox}
TCP numbers \emph{bytes}, like pages of a book posted in envelopes: each segment carries its first page number, and ``send page 1001 next'' acknowledges everything before it.
\end{intuitionbox}

\begin{definitionbox}[TCP]
A \keyword{connection-oriented, reliable} protocol with establishment, data-transfer and teardown phases; it combines GBN and SR ideas. Services: process-to-process communication, \keyword{stream delivery} (bytes through sending and receiving buffers, grouped into segments), full duplex, multiplexing/demultiplexing, connection-oriented, reliable. \src{N_Network_Systems_compressed p75-76}
\end{definitionbox}

\begin{definitionbox}[Numbering system (notes p77)]
\begin{itemize}
  \item \textbf{Byte number}: every byte sent is numbered. The first number is chosen randomly (the \keyword{ISN}, between 0 and $2^{32}-1$).
  \item \textbf{Sequence number} of a segment = the number of the \textbf{first data byte} it carries.
  \item \textbf{Acknowledgment number} = the number of the \textbf{next byte the receiver expects}; it is cumulative.
\end{itemize}
\end{definitionbox}

\textbf{Example} \src{N_Network_Systems_compressed p77}: a 5000-byte file whose first byte is numbered 10\,001, sent in five 1000-byte segments, gives segment seq numbers 10\,001, 11\,001, 12\,001, 13\,001, 14\,001 (segment 1 carries bytes 10\,001--11\,000; check: \texttt{checks/T18\_ex\_seq.py}).

\begin{table}[H]\centering\footnotesize\zebra
\begin{tabularx}{\linewidth}{L{40mm}cY}\toprule
\thead{Segment} & \thead{Seq.\ numbers consumed} & \thead{Why}\\\midrule
SYN & 1 & carries no data but must be acknowledged \\
SYN + ACK & 1 & same as SYN \\
ACK (no data) & 0 & nothing to acknowledge \\
Data segment with $n$ bytes & $n$ & one number per byte\\
FIN (no data) & 1 & must be acknowledged\\
FIN + ACK & 1 & same as FIN\\\bottomrule
\end{tabularx}
\caption{How many sequence numbers each segment consumes. \src{N_Network_Systems_compressed p82}}
\end{table}

\needspace{8\baselineskip}
\begin{question}{Worked example -- Variant of MS2024-5a}{2.5}[\tier{B}]
A TCP client's initial sequence number is 4800. It opens the connection, sends a data segment carrying 1500 bytes and then closes the
connection. Give the sequence number of (i) the SYN segment, (ii) the data segment, (iii) the FIN segment.
\end{question}
\begin{solution}[SYN consumes 1, data consumes 1500 (check: \texttt{checks/T18\_ex\_seq.py})]
(i) SYN: seq $=$ ISN $=$ \answer{4800} (consumes one number).\quad
(ii) Data: seq $=4801$ (bytes 4801--6300) $\Rightarrow$ \answer{4801}.\quad
(iii) FIN: seq $=4801+1500=$ \answer{6301}.
(The pure ACK of the handshake would also carry seq 4801 but consumes nothing.)
\end{solution}

\begin{figure}[H]\centering
\begin{tikzpicture}[x=1cm,y=1cm]
  \draw[draw=sgnavy,fill=hl!20] (0,0) rectangle (1.2,0.7); \node[font=\scriptsize\bfseries] at (0.6,0.35) {SYN};
  \draw[draw=sgnavy,fill=ok!15] (1.2,0) rectangle (10.2,0.7); \node[font=\scriptsize] at (5.7,0.35) {data: 1500 bytes};
  \draw[draw=sgnavy,fill=hl!20] (10.2,0) rectangle (11.4,0.7); \node[font=\scriptsize\bfseries] at (10.8,0.35) {FIN};
  \node[font=\scriptsize,below] at (0.6,0) {4800};
  \node[font=\scriptsize,below] at (1.75,0) {4801};
  \node[font=\scriptsize,below] at (9.65,0) {6300};
  \node[font=\scriptsize,below] at (10.8,0) {6301};
  \node[note,above] at (0.6,0.7) {1 number};
  \node[note,above] at (5.7,0.7) {1 number per byte};
  \node[note,above] at (10.8,0.7) {1 number};
\end{tikzpicture}
\caption{Sequence numbers consumed by the client in the worked example: SYN takes the ISN, data bytes 4801--6300, FIN 6301 (check: \texttt{checks/T18\_ex\_seq.py}).}
\end{figure}

\begin{warningbox}
ISN = seq of the SYN; first data byte = ISN + 1. A pure ACK consumes no number, so the next data segment reuses its seq. ACK = last byte + 1.
\end{warningbox}

\begin{answerplan}[1--2.5 marks]
One line per segment with the rule (``SYN consumes 1'') and the arithmetic. Asked in 2024 Q1(e) and Q5(a); solved on p.\,\pageref{q:MS2024-1e} and p.\,\pageref{q:MS2024-5a}.
\par\smallskip{\color{sgteal}\textbf{Self-check.}} \textbf{T18.1} A segment has seq 2001 and carries 500 bytes. What ACK number acknowledges it?\quad
\textbf{T18.2} Why is the ISN random rather than 0?
\end{answerplan}

% =====================================================================
\section{TCP segment format (header fields, flags)}\label{sec:T19}
\prob{80} \tier{B} \asked{2022 Q2(a); 2024 Q3(b); 2025 Q4(b)}

\begin{intuitionbox}
Every TCP job has a field: ports (mux), seq/ack (reliability), window (flow control), flags (connection), checksum (errors), urgent pointer.
\end{intuitionbox}

A packet in TCP is a \keyword{segment}: a \textbf{20 to 60-byte header} (20 without options) followed by data. \src{N_Network_Systems_compressed p77-78}

\begin{figure}[H]\centering
\begin{tikzpicture}[x=4.4mm,y=6.6mm]
  \foreach \b in {0,16,31}{\node[note] at ({\b+0.5},0.3) {\b};}
  \foreach \x/\w/\r/\t in {0/16/0/{Source port address (16)},16/16/0/{Destination port address (16)},
                           0/32/1/{Sequence number (32)},0/32/2/{Acknowledgment number (32)},
                           0/4/3/{HLEN\\[-2pt]\scriptsize 4 bits},4/6/3/{Reserved\\[-2pt]\scriptsize 6 bits},16/16/3/{Window size (16)},
                           0/16/4/{Checksum (16)},16/16/4/{Urgent pointer (16)}}{
    \draw[draw=sgnavy,fill=sgtint,line width=0.6pt] (\x,-\r) rectangle ++(\w,-1);
    \node[font=\small,align=center] at ({\x+\w/2},{-\r-0.5}) {\t};
  }
  \foreach \f [count=\i from 0] in {U,A,P,R,S,F}{
    \draw[draw=sgnavy,fill=hl!18,line width=0.6pt] ({10+\i},-3) rectangle ++(1,-1);
    \node[font=\scriptsize\bfseries] at ({10.5+\i},-3.5) {\f};
  }
  \draw[draw=sgnavy,fill=info!10,line width=0.6pt] (0,-5) rectangle (32,-6);
  \node[font=\small] at (16,-5.5) {Options and padding (up to 40 bytes)};
  \draw[decorate,decoration={brace,amplitude=5pt}] (32.4,0) -- node[right=6pt,note]{20 bytes\\(fixed)} (32.4,-5);
  \draw[decorate,decoration={brace,amplitude=4pt}] (32.4,-5) -- node[right=6pt,note]{0--40 B} (32.4,-6);
\end{tikzpicture}
\caption{TCP header. Flags: \textbf{U}RG, \textbf{A}CK, \textbf{P}SH, \textbf{R}ST, \textbf{S}YN, \textbf{F}IN (``UAPRSF''). \src{N_Network_Systems_compressed p78}}
\end{figure}



\needspace{16\baselineskip}
\begin{question}{Worked example -- Variant of MS2025-4b}{2}[\tier{B}]
A TCP header dump in hexadecimal is \texttt{0050D431000186A0000027115012FFFF}\dots. Find the source and destination ports, sequence number,
acknowledgment number, header length, type of segment and window size.
\end{question}
\begin{solution}[split by field width: 4 + 4 + 8 + 8 + 4 + 4 hex digits (check: \texttt{checks/T19\_ex\_tcp.py})]
\begin{tabular}{@{}lll@{}}
Source port & \texttt{0050} & \answer{80} (HTTP server)\\
Destination port & \texttt{D431} & \answer{54321}\\
Sequence number & \texttt{000186A0} & \answer{100\,000}\\
Acknowledgment number & \texttt{00002711} & \answer{10\,001}\\
HLEN & \texttt{5} & $5\times4=$ \answer{20 bytes} (no options)\\
Reserved + flags & \texttt{012} $=0000\,0001\,0010_2$ & last 6 bits $010010$: ACK and SYN set\\
Segment type & & \answer{SYN + ACK segment}\\
Window size & \texttt{FFFF} & \answer{65\,535 bytes}
\end{tabular}
\end{solution}

\needspace{8\baselineskip}
\begin{question}{Worked example -- Variant of MS2022-2a (urgent data)}{1.5}[\tier{B}]
A TCP segment of 600 bytes has HLEN $=8$, sequence number 7000, URG $=1$ and urgent pointer $=120$. How many bytes are urgent and what are their sequence numbers?
\end{question}
\begin{solution}[header first, then the pointer (check: \texttt{checks/T19\_ex\_tcp.py})]
Header $=8\times4=32$ bytes, so data $=600-32=568$ bytes, numbered 7000--7567.
Urgent pointer 120 marks where non-urgent data begins (offset 120), so urgent data $=$ \answer{120 bytes, numbers 7000--7119}; normal data 7120--7567.
\emph{Alternative:} the RFC 1122 / Forouzan reading (pointer = last urgent byte) gives 121 bytes, 7000--7120; quote it only as a remark.
\end{solution}

\begin{table}[H]\centering\footnotesize\zebra
\begin{tabularx}{\linewidth}{L{36mm}cY}\toprule
\thead{Field} & \thead{Bits} & \thead{Meaning}\\\midrule
Source / destination port & 16 + 16 & sending and receiving processes (mux/demux)\\
Sequence number & 32 & number of the first data byte in this segment (SYN: the ISN)\\
Acknowledgment number & 32 & next byte expected; valid only if ACK = 1\\
HLEN & 4 & header length in 4-byte words: 5--15 $\Rightarrow$ 20--60 bytes\\
Reserved & 6 & for future use (0)\\
Flags & 6 & URG: urgent pointer valid; ACK: acknowledgment valid; PSH: request for push (pass data up immediately); RST: reset/abort the connection; SYN: synchronise sequence numbers; FIN: terminate the connection\\
Window size & 16 & receive window \textbf{rwnd} in bytes (max 65\,535); flow control\\
Checksum & 16 & over pseudoheader (protocol 6) + header + data; \textbf{mandatory} in TCP (optional in UDP over IPv4)\\
Urgent pointer & 16 & valid if URG = 1; says where in the segment the non-urgent data begins\\
Options & 0--320 & e.g.\ MSS, window scale, timestamp\\\bottomrule
\end{tabularx}
\caption{TCP header fields to write in a ``draw and explain each field'' answer. \src{N_Network_Systems_compressed p78-79}}
\end{table}

\begin{warningbox}
HLEN is in 4-byte words (\texttt{5} = 20 B). After the HLEN digit, 3 hex digits hold 6 reserved bits + 6 flags: \texttt{002} SYN, \texttt{012} SYN+ACK, \texttt{010} ACK, \texttt{011} FIN+ACK, \texttt{018} PSH+ACK. Seq/ack are 8 hex digits each.
\end{warningbox}

\begin{answerplan}[2.5 marks, ``draw the TCP header and explain each field'']
\stepno{1} Draw the 32-bit-wide header with bit positions 0, 16, 31 and field widths (the diagram is about half the marks) \quad
\stepno{2} One line per field from the table, flags listed as UAPRSF \quad \stepno{3} Note 20--60 bytes, HLEN in words.
Asked in 2024 Q3(b), 2022 Q2(a), 2025 Q4(b); solved on p.\,\pageref{q:MS2024-3b}, p.\,\pageref{q:MS2022-2a} and p.\,\pageref{q:MS2025-4b}.
\par\smallskip{\color{sgteal}\textbf{Self-check.}} \textbf{T19.1} HLEN $=$ \texttt{F}. How many option bytes?\quad
\textbf{T19.2} Which flag is set in a segment whose flag bits are $000001$?\quad
\textbf{T19.3} Why does TCP need a window field but UDP does not?
\end{answerplan}

% =====================================================================
\section{TCP connection management: 3-way handshake, data transfer, termination}\label{sec:T20}
\prob{59} \tier{B} \asked{2022 Q3(b); 2025 Q3(a)}

\begin{intuitionbox}
Each side's starting number must be announced \emph{and} confirmed; the middle SYN+ACK does both, hence three messages.
\end{intuitionbox}

TCP is connection-oriented: it sets up a logical path in three phases, \textbf{connection establishment}, \textbf{data transfer} (both directions) and \textbf{connection termination}. \src{N_Network_Systems_compressed p80}

\textbf{Three-way handshake} (server passive open, client active open): \stepno{1} SYN (seq $=$ client ISN 8000, no data, consumes 1) \stepno{2} SYN+ACK (seq $=$ server ISN 15000, ack 8001, consumes 1, advertises rwnd) \stepno{3} ACK (seq 8001, ack 15001, consumes 0).

\stepanim{4}{2}{0.48\linewidth}{TCP connection life cycle with the numbers of the notes, one phase per frame: (1) 3-way handshake, (2) data transfer with piggybacked ACKs, (3) 3-way termination, (4) the half-close alternative. \src{N_Network_Systems_compressed p80-83}}{%
  \path[use as bounding box] (-0.95,-3.95) rectangle (7.15,0.85);
  \node[blk,minimum width=12mm,minimum height=5mm,font=\scriptsize,inner sep=1pt] at (0,0.15) {Client};
  \node[blk,minimum width=12mm,minimum height=5mm,font=\scriptsize,inner sep=1pt] at (6.2,0.15) {Server};
  \draw[lifeline] (0,-0.15) -- (0,-3.25); \draw[lifeline] (6.2,-0.15) -- (6.2,-3.25);
  \ifcase\f\or
    \node[note,anchor=east] at (-0.1,-0.4) {active};
    \node[note,anchor=west] at (6.3,-0.4) {passive};
    \draw[msg] (0,-0.4) -- node[elabel,font=\scriptsize]{SYN: seq 8000} (6.2,-0.9);
    \draw[msgok] (6.2,-1.2) -- node[elabel,font=\scriptsize]{SYN+ACK: seq 15000, ack 8001} (0,-1.7);
    \draw[msg] (0,-2.0) -- node[elabel,font=\scriptsize]{ACK: seq 8001, ack 15001} (6.2,-2.5);
    \node[note,text width=76mm] at (3.1,-3.62) {SYN and SYN+ACK consume one number each; the pure ACK none.};
  \or
    \draw[msg] (0,-0.35) -- node[elabel,font=\scriptsize]{seq 8001, ack 15001: bytes 8001--9000} (6.2,-0.8);
    \draw[msg] (0,-1.1) -- node[elabel,font=\scriptsize]{seq 9001, ack 15001: bytes 9001--10000} (6.2,-1.55);
    \draw[msgok] (6.2,-1.85) -- node[elabel,font=\scriptsize]{seq 15001, ack 10001: bytes 15001--17000} (0,-2.3);
    \draw[msg] (0,-2.6) -- node[elabel,font=\scriptsize]{ACK: seq 10001, ack 17001, rwnd 10000} (6.2,-3.05);
    \node[note,text width=76mm] at (3.1,-3.62) {ACKs ride on data (piggybacking); PSH asks for immediate delivery.};
  \or
    \draw[msg] (0,-0.4) -- node[elabel,font=\scriptsize]{FIN: seq 10001, ack 17001} (6.2,-0.9);
    \draw[msgok] (6.2,-1.2) -- node[elabel,font=\scriptsize]{FIN+ACK: seq 17001, ack 10002} (0,-1.7);
    \draw[msg] (0,-2.0) -- node[elabel,font=\scriptsize]{ACK: seq 10002, ack 17002} (6.2,-2.5);
    \node[note,text width=76mm] at (3.1,-3.62) {FIN and FIN+ACK consume one number each; the last ACK closes.};
  \or
    \draw[msg] (0,-0.35) -- node[elabel,font=\scriptsize]{FIN: seq $x$, ack $y$} (6.2,-0.75);
    \draw[msgok] (6.2,-0.95) -- node[elabel,font=\scriptsize]{ACK: seq $y-1$, ack $x+1$} (0,-1.35);
    \fill[info!12] (0.1,-1.5) rectangle (6.1,-1.95);
    \node[font=\scriptsize] at (3.1,-1.72) {server keeps sending data; client ACKs};
    \draw[msgok] (6.2,-2.1) -- node[elabel,font=\scriptsize]{FIN: seq $z$, ack $x+1$} (0,-2.5);
    \draw[msg] (0,-2.7) -- node[elabel,font=\scriptsize]{ACK: seq $x+1$, ack $z+1$} (6.2,-3.1);
    \node[note,text width=76mm] at (3.1,-3.62) {Half-close: one direction closes first, the other keeps sending.};
  \fi
}

\textbf{Data transfer}: ACKs are \keyword{piggybacked} on data; PSH asks for immediate delivery; every segment advertises rwnd. \textbf{Termination}: 3-way (FIN, FIN+ACK, ACK) or \keyword{half-close} (frame 4); FIN consumes one number. \src{N_Network_Systems_compressed p81-83}



\needspace{8\baselineskip}
\begin{question}{Worked example -- Variant of MS2022-3b}{2.5}[\tier{C}]
MSS $=500$ bytes both ways. Client ISN $=200$, server ISN $=500$. The client sends 1000 bytes, the server sends 500 bytes, no loss,
the application reads data at once. List every segment with its SYN/ACK/FIN bits, SEQ and ACK numbers.
\end{question}
\begin{solution}[trace (check: \texttt{checks/T20\_ex\_trace.py})]
\centering\small
\begin{tabular}{@{}clcccrr@{}}\toprule
\# & Direction, content & SYN & ACK & FIN & SEQ & ACK no.\\\midrule
1 & C$\to$S SYN & 1 & 0 & 0 & 200 & --\\
2 & S$\to$C SYN+ACK & 1 & 1 & 0 & 500 & 201\\
3 & C$\to$S ACK & 0 & 1 & 0 & 201 & 501\\
4 & C$\to$S data 201--700 & 0 & 1 & 0 & 201 & 501\\
5 & C$\to$S data 701--1200 & 0 & 1 & 0 & 701 & 501\\
6 & S$\to$C data 501--1000 (+ ACK) & 0 & 1 & 0 & 501 & 1201\\
7 & C$\to$S FIN (+ ACK) & 0 & 1 & 1 & 1201 & 1001\\
8 & S$\to$C FIN+ACK & 0 & 1 & 1 & 1001 & 1202\\
9 & C$\to$S ACK & 0 & 1 & 0 & 1202 & 1002\\\bottomrule
\end{tabular}\par\raggedright
Assumptions: the server's data ACK is piggybacked (row 6) and the client closes first with a 3-way close; a separate pure ACK after row 5 would not change any number.
\end{solution}

\begin{warningbox}
The handshake ACK and the first data segment carry the same seq (8001). After a FIN the peer acknowledges FIN seq + 1. In the half-close the server's pure ACK has seq $y-1$ (it consumes nothing).
\end{warningbox}

\begin{answerplan}[3 marks, ``explain establishment, data transfer and termination'']
\stepno{1} Name the three phases \quad \stepno{2} Draw the handshake with SYN / SYN+ACK / ACK and their seq/ack numbers, active vs passive open \quad
\stepno{3} One data segment pair showing piggybacking \quad \stepno{4} FIN / FIN+ACK / ACK (mention half-close) \quad
Keywords: \keyword{ISN}, \keyword{piggybacking}, \keyword{PSH}, \keyword{half-close}. Asked in 2025 Q3(a) and 2022 Q3(b); solved on p.\,\pageref{q:MS2025-3a} and p.\,\pageref{q:MS2022-3b}.
\par\smallskip{\color{sgteal}\textbf{Self-check.}} \textbf{T20.1} Why is a 2-way handshake not enough?\quad
\textbf{T20.2} In the storyboard, what seq would a further 100-byte client data segment carry if sent before the FIN?
\end{answerplan}

% =====================================================================
\section{TCP state transition diagram}\label{sec:T21}
\prob{35} \tier{C} \asked{2024 Q5(b)}

\begin{intuitionbox}
The state diagram is \S\ref{sec:T20} drawn as a map: every arrow is ``event / action''; client path on the left, server path on the right.
\end{intuitionbox}

\begin{definitionbox}[TCP finite state machine]
TCP keeps a state for each connection. A transition is written \textsf{input / output}: for example \textsf{SYN / SYN+ACK} means ``on receiving a SYN, send a SYN+ACK'';
``--'' means nothing is sent. The diagram has 11 states. \src{N_Network_Systems_compressed p84-85}
\end{definitionbox}

\begin{figure}[!htbp]\centering
\begin{tikzpicture}[x=1cm,y=0.74cm,
   st/.style={smstate,minimum width=21mm,minimum height=6.5mm,font=\scriptsize\sffamily},
   cli/.style={->,line width=0.9pt,color=sgnavy},
   srv/.style={->,line width=0.9pt,color=ok,dashed},
   bth/.style={->,line width=0.9pt,color=info},
   lb/.style={elabel,font=\scriptsize}]
  \node[st] (CL) at (4.5,0) {CLOSED};
  \node[st] (LI) at (4.5,-1.8) {LISTEN};
  \node[st] (SR) at (0,-3.6) {SYN-RCVD};
  \node[st] (SS) at (9,-3.6) {SYN-SENT};
  \node[st,minimum width=26mm,fill=hl!15,draw=hl] (ES) at (4.5,-5.4) {ESTABLISHED};
  \node[st] (F1) at (0,-7.4) {FIN-WAIT-1};
  \node[st] (CG) at (4.5,-7.4) {CLOSING};
  \node[st] (CW) at (9,-7.4) {CLOSE-WAIT};
  \node[st] (F2) at (0,-9.6) {FIN-WAIT-2};
  \node[st] (TW) at (4.5,-9.6) {TIME-WAIT};
  \node[st] (LA) at (9,-9.6) {LAST-ACK};
  \draw[arr] (4.5,0.9) node[above,note]{start} -- (CL.north);
  \draw[srv] (CL) -- node[lb]{passive open / --} (LI);
  \draw[cli] (CL.east) -| node[lb,pos=0.3]{active open / SYN} (SS.north);
  \draw[srv] (LI.west) -| node[lb,pos=0.3]{SYN / SYN+ACK} (SR.north);
  \draw[bth] (SS.west) -- node[lb]{SYN / SYN+ACK (simultaneous open)} (SR.east);
  \draw[srv] (SR.south) |- node[lb,pos=0.3]{ACK / --} (ES.west);
  \draw[cli] (SS.south) |- node[lb,pos=0.3]{SYN+ACK / ACK} (ES.east);
  \draw[cli] (ES.south west) -- node[lb]{close / FIN} (F1.north);
  \draw[srv] (ES.south east) -- node[lb]{FIN / ACK} (CW.north);
  \draw[cli] (F1) -- node[lb]{ACK / --} (F2);
  \draw[bth] (F1) -- node[lb]{FIN / ACK} (CG);
  \draw[bth] (F1.south east) -- node[lb,pos=0.45]{FIN+ACK / ACK} (TW.north west);
  \draw[bth] (CG) -- node[lb]{ACK / --} (TW);
  \draw[cli] (F2) -- node[lb]{FIN / ACK} (TW);
  \draw[srv] (CW) -- node[lb]{close / FIN} (LA);
  \draw[cli] (TW.south) -- (4.5,-10.5) -- node[lb]{timeout (2MSL) / --} (-1.9,-10.5) -- (-1.9,0) -- (CL.west);
  \draw[srv] (LA.east) -- (10.9,-9.6) -- node[lb]{ACK / --} (10.9,0.55) -- (4.9,0.55) -- ([xshift=4mm]CL.north);
  \begin{scope}[shift={(0,-11.3)}]
    \draw[cli] (0,0) -- (1,0); \node[note,anchor=west] at (1.1,0) {client transition};
    \draw[srv] (3.8,0) -- (4.8,0); \node[note,anchor=west] at (4.9,0) {server transition};
    \draw[bth] (7.6,0) -- (8.6,0); \node[note,anchor=west] at (8.7,0) {client or server};
  \end{scope}
\end{tikzpicture}
\caption{TCP state transition diagram (labels are \textsf{input / output}). \src{N_Network_Systems_compressed p84}}
\end{figure}



\begin{table}[!htbp]\centering\footnotesize\zebra
\begin{tabularx}{\linewidth}{L{22mm}YL{20mm}Y}\toprule
\thead{State} & \thead{Meaning} & \thead{State} & \thead{Meaning}\\\midrule
CLOSED & no connection exists & FIN-WAIT-2 & ACK of first FIN received; waiting for second FIN\\
LISTEN & passive open; waiting for SYN & CLOSE-WAIT & first FIN received, ACK sent; waiting for application to close\\
SYN-SENT & SYN sent; waiting for ACK & TIME-WAIT & second FIN received, ACK sent; waiting 2MSL\\
SYN-RCVD & SYN+ACK sent; waiting for ACK & LAST-ACK & second FIN sent; waiting for ACK\\
ESTABLISHED & data transfer in progress & CLOSING & both sides closing simultaneously\\
FIN-WAIT-1 & first FIN sent; waiting for ACK &  & \\\bottomrule
\end{tabularx}
\caption{The 11 TCP states. \src{N_Network_Systems_compressed p85}}
\end{table}

\textbf{Common scenario} \src{N_Network_Systems_compressed p85-86}. \emph{Client:} CLOSED $\to$ SYN-SENT $\to$ ESTABLISHED $\to$ FIN-WAIT-1 $\to$ FIN-WAIT-2 $\to$ TIME-WAIT $\to$ CLOSED. \emph{Server:} CLOSED $\to$ LISTEN $\to$ SYN-RCVD $\to$ ESTABLISHED $\to$ CLOSE-WAIT $\to$ LAST-ACK $\to$ CLOSED.

\textbf{Why TIME-WAIT lasts 2MSL} (MSL = maximum segment lifetime): the client can resend the final ACK if it was lost, and old duplicates of this connection die out before the socket pair is reused.

\begin{warningbox}
CLOSE-WAIT is on the passive side; FIN-WAIT-1/2 and TIME-WAIT on the active closer's side. Label every arrow ``event / action''.
\end{warningbox}

\begin{answerplan}[2.5 marks, ``mention the TCP states and draw the state transition diagram'']
\stepno{1} List the 11 states with one-line meanings (table) \quad \stepno{2} Draw the FSM with client (solid) and server (dashed) paths and labelled arrows \quad
\stepno{3} One line on TIME-WAIT (2MSL). Asked in 2024 Q5(b); solved on p.\,\pageref{q:MS2024-5b}.
\par\smallskip{\color{sgteal}\textbf{Self-check.}} \textbf{T21.1} In which state does a server wait for its own application to close?\quad
\textbf{T21.2} Which state is entered when both sides send FIN at the same time?
\end{answerplan}

% =====================================================================
\section{Windows in TCP and TCP flow control}\label{sec:T22}
\prob{41} \tier{C} \asked{2025 Q1(d)}

\begin{intuitionbox}
The receiver keeps advertising its free buffer space (\keyword{rwnd}); the sender never has more unacknowledged bytes in flight than that, or than cwnd.
\end{intuitionbox}

\begin{definitionbox}[TCP windows]
\textbf{Send window} (bytes, variable) $=\min(rwnd,\ cwnd)$. \textbf{Receive window} $rwnd=$ buffer size $-$ bytes waiting to be pulled by the process, advertised in the window field of every segment (the notes' data-transfer figure shows rwnd $=10\,000$). \src{N_Network_Systems_compressed p81, p87}
\end{definitionbox}

\begin{figure}[H]\centering
\begin{tikzpicture}[x=1cm,y=1cm]
  \foreach \a/\b/\f/\t in {0/2.4/black!6/{sent, ACKed},2.4/6.4/hl!25/{outstanding (no ACK yet)},6.4/9.4/ok!15/{can be sent now},9.4/12.4/white/{cannot be sent}}{
    \draw[draw=sgslate,fill=\f] (\a,0) rectangle (\b,0.7); \node[font=\scriptsize] at ({(\a+\b)/2},0.35) {\t};}
  \draw[hl,line width=1.3pt] (2.4,0) rectangle (9.4,0.7);
  \draw[arr] (2.4,1.4) node[above,font=\scriptsize]{$S_f$} -- (2.4,0.75);
  \draw[arr] (6.4,1.4) node[above,font=\scriptsize]{$S_n$} -- (6.4,0.75);
  \draw[<->,sgslate] (2.4,-0.3) -- node[elabel]{send window $=\min(rwnd,cwnd)$} (9.4,-0.3);
  \node[font=\scriptsize\bfseries,anchor=east] at (-0.1,0.35) {Sender};
  \begin{scope}[yshift=-2.3cm]
  \foreach \a/\b/\f/\t in {0/2.4/black!6/{read by process},2.4/6.4/info!20/{received, waiting: 2500 B},6.4/10.0/ok!15/{rwnd = free: 3500 B},10.0/12.4/white/{cannot receive}}{
    \draw[draw=sgslate,fill=\f] (\a,0) rectangle (\b,0.7); \node[font=\scriptsize] at ({(\a+\b)/2},0.35) {\t};}
  \draw[decorate,decoration={brace,mirror,amplitude=5pt}] (2.4,-0.1) -- node[below=5pt,note]{receive buffer = 6000 B} (10.0,-0.1);
  \node[font=\scriptsize\bfseries,anchor=east] at (-0.1,0.35) {Receiver};
  \end{scope}
\end{tikzpicture}
\caption{TCP send window (bytes) and receive window. Numbers match the worked example below.}
\end{figure}

The window \textbf{opens} when the receiver advertises more space, \textbf{closes} as bytes are acknowledged and should not \textbf{shrink}; rwnd $=0$ stops the sender, which probes with a persist timer. Unlike generic push/pull flow control (\S\ref{sec:T13}), TCP's is \keyword{byte-oriented}, \keyword{receiver-advertised} in every segment and also capped by cwnd.

\needspace{8\baselineskip}
\begin{question}{Worked example -- Expected (cluster C35)}{1}[\tier{C}]
A receiver has a 6000-byte buffer; 2500 bytes have been received but not yet read by the process. The sender's cwnd is 5 MSS with MSS $=1000$ bytes. Find rwnd and the sender's window.
\end{question}
\begin{solution}[(check: \texttt{checks/T22\_ex\_rwnd.py})]
$rwnd=6000-2500=\answer{3500 B}$; send window $=\min(3500,\ 5\times1000)=\answer{3500 B}$ (flow control, not congestion, is the limit here).
\end{solution}

\textbf{Window size vs MSS} (2025 Q1d): the \keyword{window} is the number of bytes that may be outstanding without an ACK, set by the receiver (rwnd, 16-bit window field) and the congestion state (cwnd), and it changes with every segment. The \keyword{MSS} is the largest amount of \emph{data} in one segment, negotiated once as a TCP option in the SYN segments and fixed for the connection. A window usually holds several MSS-sized segments.

\begin{warningbox}
The window field limits \emph{outstanding bytes}, not the size of one segment: ``window = MSS'' scores zero.
\end{warningbox}

\begin{answerplan}[1 mark]
Two definitions in one line each, plus ``window changes dynamically / MSS is fixed''. Asked in 2025 Q1(d); solved on p.\,\pageref{q:MS2025-1d}.
\par\smallskip{\color{sgteal}\textbf{Self-check.}} \textbf{T22.1} Buffer 4000 B, 1500 B waiting to be read: rwnd?\quad \textbf{T22.2} What must the sender do when rwnd $=0$?
\end{answerplan}

% =====================================================================
\section{TCP congestion control: slow start, congestion avoidance, fast recovery; throughput}\label{sec:T23}
\prob{59} \tier{B} \asked{2022 Q4(c); 2025 Q5(iii)}

\begin{intuitionbox}
Probe the network: double per RTT up to a threshold, then add one segment per RTT; on loss, cut back hard. The result is the \keyword{AIMD} sawtooth.
\end{intuitionbox}

The notes list three ingredients, a \keyword{congestion window} (cwnd), \keyword{congestion detection} and \keyword{congestion policies} (slow start, congestion avoidance), built on two mechanisms: additive increase and multiplicative decrease. \src{N_Network_Systems_compressed p87}

\begin{itemize}
  \item \textbf{Slow start} (exponential): cwnd starts at 1 MSS, $+1$ MSS per ACK, so it doubles every RTT ($1,2,4,8,\ldots$) up to \keyword{ssthresh}. \textbf{Congestion avoidance} (additive): once cwnd $\ge$ ssthresh, $+1$ MSS per RTT. \src{N_Network_Systems_compressed p87-89}
  \item \textbf{Detection}: timeout = strong congestion; 3 duplicate ACKs = mild (one segment lost). \textbf{Multiplicative decrease}: ssthresh $=$ cwnd/2.
  \item \textbf{Tahoe}: either signal $\Rightarrow$ cwnd $=1$, slow start. \textbf{Reno}: timeout the same; 3 dup ACKs $\Rightarrow$ \keyword{fast recovery}: cwnd $=$ ssthresh $+3$, $+1$ per further dup ACK; on the new ACK cwnd $=$ ssthresh and congestion avoidance resumes. (Simplified Reno skips the $+3$; the value after fast recovery, cwnd $=$ ssthresh, is the same.)
\end{itemize}

\begin{figure}[!htbp]\centering
\begin{tikzpicture}
\begin{axis}[width=0.92\linewidth,height=4.6cm,xlabel={Transmission round (RTT)},ylabel={cwnd (MSS)},
  grid=major,grid style={sggrid},xmin=0,xmax=20,ymin=0,ymax=23,xtick={1,3,...,19},ytick={0,4,...,20},
  every axis plot/.append style={line width=1pt},legend style={font=\scriptsize,at={(0.99,0.99)},anchor=north east}]
  \addplot[sgnavy,mark=*,mark size=1.6pt] coordinates {(1,1)(2,2)(3,4)(4,8)(5,16)(6,17)(7,18)(8,19)(9,20)(10,1)(11,2)(12,4)(13,8)(14,10)(15,11)(16,12)(17,6)(18,7)(19,8)};
  \addlegendentry{cwnd}
  \addplot[sgred,dashed,const plot] coordinates {(0.5,16)(9.5,10)(16.5,6)(19.5,6)};
  \addlegendentry{ssthresh}
  \node[elabel,text=sgred] at (axis cs:9,21.8) {timeout};
  \node[elabel,text=sgred] at (axis cs:16,13.6) {3 dup ACKs};
  \node[elabel] at (axis cs:3.0,9.5) {slow start};
  \node[elabel] at (axis cs:7.0,13.6) {cong.\ avoidance};
  \node[elabel] at (axis cs:18.0,3.2) {fast recovery $\to$ CA};
\end{axis}
\end{tikzpicture}
\caption{cwnd trace of the worked example (Reno): slow start to ssthresh 16, additive increase, timeout at 20 (ssthresh 10, cwnd 1), three duplicate ACKs at 12 (ssthresh 6, cwnd 6).}
\end{figure}





\needspace{8\baselineskip}
\begin{question}{Worked example -- Expected (cluster C36)}{2.5}[\tier{B}]
TCP Reno, cwnd starts at 1 MSS, ssthresh $=16$ MSS. A timeout occurs in round 9 and three duplicate ACKs arrive in round 16. Give cwnd for rounds 1--19.
\end{question}
\begin{solution}[apply the three rules round by round (check: \texttt{checks/T23\_ex\_cwnd.py})]
Rounds 1--5 slow start: 1, 2, 4, 8, 16 (reaches ssthresh). Rounds 6--9 additive increase: 17, 18, 19, 20.
Timeout at 20 $\Rightarrow$ ssthresh $=10$, cwnd $=1$. Rounds 10--14: 1, 2, 4, 8, 10 (capped at ssthresh); 15--16: 11, 12.
3 dup ACKs at 12 $\Rightarrow$ ssthresh $=6$, cwnd $=6+3=9$ during fast recovery, and cwnd $=$ ssthresh $=6$ when the new ACK arrives. Rounds 17--19: \answer{6, 7, 8}.
Assumption: one cwnd value per round, slow start capped at ssthresh. Under Tahoe, round 17 would restart at cwnd $=1$.
\end{solution}

\begin{figure}[!htbp]\centering
\begin{tikzpicture}[x=1cm,y=1cm,st/.style={smstate,minimum width=30mm,minimum height=9mm,font=\scriptsize\sffamily},lb/.style={elabel,font=\scriptsize,align=center}]
  \node[st] (SS) at (0,0) {Slow start};
  \node[st] (CA) at (8,0) {Congestion avoidance};
  \node[st,fill=hl!15,draw=hl] (FR) at (4,-3.8) {Fast recovery (Reno)};
  \draw[arr] (SS) -- node[lb]{cwnd $\ge$ ssthresh} (CA);
  \draw[arrbad] (CA.north) -- (8,1.2) -- node[lb]{timeout: ssthresh $=$ cwnd/2, cwnd $=1$} (0,1.2) -- (SS.north);
  \draw[arrbad] (FR.west) -| node[lb,pos=0.7]{timeout:\\cwnd $=1$} ([xshift=-8mm]SS.south);
  \draw[arrbad] (SS.south east) -- node[lb]{3 dup ACKs:\\ssthresh $=$ cwnd/2\\cwnd $=$ ssthresh $+3$} (FR.north west);
  \draw[arrbad] (CA.south west) -- node[lb]{3 dup ACKs:\\ssthresh $=$ cwnd/2\\cwnd $=$ ssthresh $+3$} (FR.north east);
  \draw[arrok] (FR.east) -| node[lb,pos=0.75]{new ACK:\\cwnd $=$ ssthresh} ([xshift=8mm]CA.south);
  \draw[arr] (SS.south west) .. controls +(-1.3,-0.2) and +(-1.3,0.2) .. node[lb,left]{ACK:\\cwnd $+1$} (SS.north west);
  \draw[arr] (CA.south east) .. controls +(1.3,-0.2) and +(1.3,0.2) .. node[lb,right]{ACK:\\cwnd $+1/$cwnd} (CA.north east);
  \draw[arr] (FR.south west) .. controls +(-0.5,-0.9) and +(0.5,-0.9) .. node[lb,below]{dup ACK: cwnd $+1$} (FR.south east);
\end{tikzpicture}
\caption{Congestion-control policy transitions (Reno, per conventions). Tahoe has no fast-recovery state: every detection sends it back to slow start with cwnd $=1$.}
\end{figure}

\begin{formulabox}[TCP throughput (notes p90)]
$\text{Throughput}=\dfrac{0.75\,W_{max}}{RTT}$, where $W_{max}$ = average window (bytes or bits) at the moments congestion is detected (sawtooth peaks); the window swings between $W_{max}/2$ and $W_{max}$, so its mean is $0.75\,W_{max}$. \src{N_Network_Systems_compressed p90}
\end{formulabox}

\needspace{8\baselineskip}
\begin{question}{Worked example -- from the notes}{2.5}[\tier{D}]
MSS $=10$ KB and RTT $=100$ ms. The cwnd sawtooth peaks at 10, 12, 10, 8 and 8 MSS. Find the throughput. \src{N_Network_Systems_compressed p90}
\end{question}
\begin{solution}[average peak, then the formula (check: \texttt{checks/T23\_ex\_throughput.py})]
$W_{max}=\frac{10+12+10+8+8}{5}=9.6$ MSS $=96$ KB $=96\times1024\times8=786\,432$ bits (1 KB $=1024$ B).\\
Throughput $=\frac{0.75\times786\,432}{0.1}=\answer{5.90 Mbps}$ (720 KB/s). With 1 KB $=1000$ B: 5.76 Mbps.
\end{solution}

\begin{examtip}[the notes' figure]
The notes write $0.75\times960$ kbps$/100$ ms $=7.2$ Mbps, treating 96 KB as 960 kilobits; converting KB to bits ($\times8$) gives 5.90 Mbps. Show the conversion line.
\end{examtip}

\begin{table}[H]\centering\footnotesize\zebra
\begin{tabularx}{\linewidth}{L{24mm}YY}\toprule
\thead{Aspect} & \thead{Slow start} & \thead{Additive increase (congestion avoidance)}\\\midrule
Growth & exponential: doubles every RTT & linear: $+1$ MSS every RTT\\
Per ACK & cwnd $+=1$ MSS & cwnd $+=1/\text{cwnd}$ MSS\\
When & connection start, after a timeout, while cwnd $<$ ssthresh & once cwnd $\ge$ ssthresh\\
Sequence & $1,2,4,8,\ldots$ & $i,\ i+1,\ i+2,\ldots$\\\bottomrule
\end{tabularx}
\caption{Slow start vs additive increase (2022 Q4c short note).}
\end{table}

\begin{warningbox}
``Slow'' start is the \emph{exponential} phase. Timeout $\Rightarrow$ cwnd $=1$ in Tahoe and Reno; only 3 duplicate ACKs differ. ssthresh $=$ half the cwnd \emph{at detection}.
\end{warningbox}

\begin{answerplan}[2.5 marks, short note ``TCP congestion control'']
\stepno{1} Define cwnd and ssthresh; send window $=\min(rwnd,cwnd)$ \quad \stepno{2} Slow start, congestion avoidance, detection (timeout vs 3 dup ACKs), fast recovery, each one line \quad
\stepno{3} Draw the cwnd-vs-RTT sawtooth with ssthresh marked \quad Keywords: \keyword{AIMD}, \keyword{Tahoe}, \keyword{Reno}. Asked in 2025 Q5(iii) and 2022 Q4(c); solved on p.\,\pageref{q:MS2025-5iii} and p.\,\pageref{q:MS2022-4c}.
\par\smallskip{\color{sgteal}\textbf{Self-check.}} \textbf{T23.1} ssthresh $=8$, cwnd starts at 1, no loss: cwnd in rounds 1--6?\quad
\textbf{T23.2} Reno, cwnd $=24$ when 3 duplicate ACKs arrive: new ssthresh and cwnd?\quad
\textbf{T23.3} $W_{max}=20$ MSS, MSS $=1000$ B, RTT $=50$ ms: throughput?
\end{answerplan}

% =====================================================================
\section{Quality of Service (QoS)}\label{sec:T24}
\prob{8} \tier{D} \asked{not yet asked; listed in the lesson plan, not taught in the notes}

\begin{intuitionbox}
A file transfer wants every byte; a voice call wants bytes on time. QoS is the set of characteristics a flow needs and the techniques that provide them.
\end{intuitionbox}

\begin{definitionbox}[Quality of Service]
The ability of a network to give a \keyword{flow} (packets from a source to a destination) the service it needs, measured by reliability, delay, jitter and bandwidth. Named in the lesson plan but not taught in the notes, so this follows the textbook. \src{SYL_Lesson_Plan_NS_sample p2}
\end{definitionbox}

\begin{table}[H]\centering\footnotesize\zebra
\begin{tabularx}{\linewidth}{L{30mm}Y}\toprule
\thead{Item} & \thead{Meaning}\\\midrule
Flow characteristics & \textbf{reliability} (losses force retransmission: file transfer, e-mail); \textbf{delay} (telephony, conferencing); \textbf{jitter}, the variation in delay (live audio/video); \textbf{bandwidth} (bits/s the flow needs)\\
Scheduling & FIFO queuing, priority queuing, weighted fair queuing\\
Traffic shaping & \textbf{leaky bucket}: constant output rate, bursts smoothed; \textbf{token bucket}: tokens saved while idle allow limited bursts\\
Reservation, admission & reserve buffers/bandwidth in advance (RSVP); a router admits or rejects a flow by its specification and current load\\
Service models & Integrated Services (per-flow, RSVP) vs Differentiated Services (per-class marking, scalable)\\\bottomrule
\end{tabularx}
\caption{QoS characteristics and techniques.}
\end{table}

\begin{answerplan}[2.5 marks, if it appears]
Define QoS and flow; list the four characteristics with one application each; name four techniques, explaining leaky vs token bucket in one line each.
\par\smallskip{\color{sgteal}\textbf{Self-check.}} \textbf{T24.1} Which QoS characteristic matters most for a live voice call, and why?\quad \textbf{T24.2} How does a token bucket differ from a leaky bucket?
\end{answerplan}

% =====================================================================
\section*{Unit recap}
\begin{itemize}
  \item Process-to-process delivery; socket = IP + port; ports 0--1023 well-known, 1024--49151 registered, 49152--65535 dynamic.
  \item $U_{SW}=\frac{1}{1+2a}$, $a=\frac{T_p}{T_t}$; window $W$: $U=\min(1,\frac{W}{1+2a})$; BDP $=B\times RTT$. ``Every $k$-th lost'' counts transmissions: SW $T=N+\lfloor\frac{N-1}{k-1}\rfloor$, GBN: draw the strip.
  \item SW window 1, ACK = next expected. GBN $S_{size}\le2^m-1$, cumulative ACK, resend all outstanding. SR $S_{size}=R_{size}=2^{m-1}$, individual ACK, resend only lost.
  \item UDP 8-byte header (src, dst, total length, checksum; 0 = not computed). DNS: UDP normally, TCP for $>512$ B and zone transfer.
  \item TCP numbers bytes: seq = first data byte, ack = next byte expected; SYN, SYN+ACK, FIN consume 1, pure ACK 0.
  \item TCP header 20--60 B, HLEN in words, flags UAPRSF, window = rwnd, urgent bytes seq $\ldots$ seq+UP$-1$.
  \item Handshake SYN $x$ / SYN+ACK ($y$, ack $x+1$) / ACK; close FIN, FIN+ACK, ACK or half-close; 11 states, TIME-WAIT = 2MSL.
  \item Window $=\min(rwnd,cwnd)$; slow start doubles to ssthresh, then $+1$/RTT; detection: ssthresh $=$ cwnd/2; timeout: cwnd $=1$; Reno 3 dup ACKs: fast recovery, then cwnd $=$ ssthresh; throughput $=0.75W_{max}/RTT$.
\end{itemize}

\needspace{12\baselineskip}
\section*{Self-check answers}
\begin{multicols}{2}\footnotesize
\textbf{T12.1} Dynamic/private (49152--65535): an ephemeral port chosen by a client. \textbf{T12.2} The destination port number.

\textbf{T13.1} The consumer asks only when it is ready, so it can never be overwhelmed. \textbf{T13.2} Detect and discard corrupted packets; resend lost or discarded ones; discard duplicates; buffer out-of-order packets. \textbf{T13.3} It assumes no loss, duplication or reordering, so nothing has to be identified.

\textbf{T14.1} When an ACK is lost the packet is resent; the number lets the receiver recognise the duplicate. \textbf{T14.2} $T_t=1$ ms, $T_p=10$ ms, $a=10$, $U=1/21=4.76\%$ (short form 5\%). \textbf{T14.3} $8+\lfloor7/4\rfloor=9$.

\textbf{T15.1} $2^4-1=15$. \textbf{T15.2} 11: 1 2 3 [4 lost] 5 6 $|$ 4 [5 lost] 6 $|$ 5 6. \textbf{T15.3} Packets up to 4 arrived; the receiver expects 5; $S_f$ moves to 5.

\textbf{T16.1} $2^m-1\ge31\Rightarrow m=5$. \textbf{T16.2} SR resends a single packet when its own timer expires; GBN always resends from the oldest outstanding packet, so one timer on $S_f$ is enough.

\textbf{T17.1} $64-8=56$ bytes. \textbf{T17.2} 17. \textbf{T17.3} A late retransmission is useless for real-time playout; no connection set-up and no congestion throttling keep delay low (UDP also supports multicast).

\textbf{T18.1} 2501. \textbf{T18.2} So segments of an old connection are not mistaken for a new one (and sequence numbers are hard to guess).

\textbf{T19.1} $15\times4=60$ B header $\Rightarrow$ 40 option bytes. \textbf{T19.2} FIN. \textbf{T19.3} TCP does flow control (rwnd advertised in the window field); UDP has none.

\textbf{T20.1} The server's ISN would never be confirmed, and an old duplicate SYN could open a false connection; the third segment confirms the server's number. \textbf{T20.2} seq 10001 (the FIN would then carry 10101).

\textbf{T21.1} CLOSE-WAIT. \textbf{T21.2} CLOSING.

\textbf{T22.1} $4000-1500=2500$ B. \textbf{T22.2} Stop sending new data and probe (persist timer) until a non-zero rwnd is advertised.

\textbf{T23.1} 1, 2, 4, 8, 9, 10. \textbf{T23.2} ssthresh $=12$; cwnd $=12+3=15$ during fast recovery, 12 after the new ACK. \textbf{T23.3} $0.75\times20\times8000/0.05=2.4$ Mbps.

\textbf{T24.1} Delay (and jitter): late speech is useless, a lost sample is tolerable. \textbf{T24.2} A leaky bucket sends at a fixed rate regardless of idle time; a token bucket saves tokens while idle and so allows controlled bursts.
\end{multicols}
